\begin{frame}{keine Zeit mehr lol}
    \begin{exercise}
        \only<1>{Berechnen Sie mithilfe des Horner-Schemas die Hexadezimaldarstellung von \(1494,59375\).}
        \only<2>{Wandeln Sie das Ergebnis in das Oktalsystem um, indem Sie die Stellen in geeignete Ziffernfolgen zerlegen und anschließend zusammenfassen.}
        \only<3>{Überprüfen Sie Ihre Rechnung, indem Sie die Zahl in das Dezimalsystem umwandeln.}
    \end{exercise}
    \begin{solve}
        \only<1>{\(1494,59375_{10} = 5\text{D}6,98_{16}\)}%
        \only<2>{\faAngleRight~Tafel}%
        \only<3>{\begin{align*}
            2726,46_{8} &= 2\cdot8^3 + 7\cdot8^2 + 2\cdot8^1 + 6\cdot8^0 + 4\cdot 8^{-1} + 6\cdot8^{-2}\\
            &= 1024 + 448 + 16 + 6 + 0,5 + 0,09375\\
            &= 1494,59375_{10}
        \end{align*}}
    \end{solve}
\end{frame}